\section{Architecture}
\label{sec:architecture}

\begin{outline}
  \item Three layers per game, each in its own place and each optional until needed: the game's own code, pinned (a package or a checked-out repository); a thin environment package when the game ships no Gymnasium environment; and one adapter per backend, which builds the environment, maps controls to actions, renders frames and sound, and reports what to record.
  \item The run never knows the engine: it calls only the adapter; game engines are imported only by adapters; the layers import only downwards, and a contract checked in continuous integration enforces it.
  \item Three kinds of settings, kept apart because they are shared differently: the curriculum (what is played, the same for every participant and site), the rig (how this setup plays every run), and the run's own numbers (participant, session, run, duration). Neither file names the other; what needs both is decided when the run starts.
  \item One run is one process, one command and one data folder; a session is a plain shell script with one line per run, written by hand or by a tool, and managed, if wanted, by a session launcher \TODO{name the compatible launchers}.
  \item Tools that can run on their own stay outside the run: making and measuring a rig, calibrating a controller, recording the session's input devices, converting logs. The run only reads what they filed, at start-up.
  \item Checks at the desk: every problem that can be found before the participant's screen opens stops the run there, with the fix in the message; a default that changes what is recorded is printed and recorded, never silent.
  \item Install only what is needed: each backend is an optional extra, imported only when a curriculum names it.
  \item Table: the 15 backends, their action spaces, how their savestates work, and their platforms \TODO{generate from the software's backend list}.
\end{outline}
